\documentclass{article}

\usepackage[margin=1.2cm, top=1cm, bottom=1.8cm]{geometry}
\pagestyle{plain}
\linespread{1.4}

\usepackage{fontspec}
\usepackage{unicode-math}
\usepackage[main=greek, english]{babel}

\setmainfont{CMU Serif}
\setsansfont{CMU Sans Serif}
\setmonofont{CMU Typewriter Text}
\setmathfont{Latin Modern Math}

\usepackage{amsmath, enumitem}
\usepackage{dirtytalk, caption, subcaption, array, graphicx}
\usepackage{titlesec}
\usepackage{soul}
\usepackage[unicode]{hyperref}
\usepackage{verbatim}

\usepackage{pdfpages}

\usepackage{tikz}
\usetikzlibrary{arrows.meta, positioning}

\hypersetup{
    hidelinks,
    bookmarksopen=true,
    bookmarksnumbered=true,
    pdftitle={Συστήματα Αναμονής - 3η Εργαστηριακή Άσκηση},
    pdfauthor={Πέππας Μιχαήλ-Αθανάσιος}
}

\newcolumntype{P}[1]{>{\centering\arraybackslash}p{#1}}
\titleformat*{\subsubsection}{\large\bfseries}

\begin{document}

\begin{center}
    {
    \Large
    {
    \textbf{ΕΘΝΙΚΟ ΜΕΤΣΟΒΙΟ ΠΟΛΥΤΕΧΝΕΙΟ\\
    ΣΧΟΛΗ ΗΛΕΚΤΡΟΛΟΓΩΝ ΜΗΧΑΝΙΚΩΝ\\
    ΚΑΙ ΜΗΧΑΝΙΚΩΝ ΥΠΟΛΟΓΙΣΤΩΝ\\}
    }

    \vspace{1.5cm}
    \rule{18cm}{0.5pt}

    Τομέας Επικοινωνιών, Ηλεκτρονικής και Συστημάτων Πληροφορικής\\
    Εργαστήριο Διαχείρισης και Βέλτιστου Σχεδιασμού Δικτύων Τηλεματικής - NETMODE\\

    \rule{18cm}{0.5pt}

    \textbf{Συστήματα Αναμονής}\\
    \textbf{$\mathbf{3^{η}}$ Εργαστηριακή Άσκηση}\\
    }
\end{center}

\vspace{0.8cm}
\begin{center}
\includegraphics[width=7cm]{EMP Logo.png}
\end{center}

\Large
{
\vspace{0.8cm}
\noindent
\textbf{Στοιχεία Φοιτητή}
\begin{itemize}
    \item Όνομα: Πέππας Μιχαήλ-Αθανάσιος
    \item Αριθμός Μητρώου: 03121026
    \item $3^{\eta}$ Εργαστηριακή Ομάδα (Δευτέρα, 12:45-14:30)
    \item Ημερομηνία Διεξαγωγής Εργαστηρίου: 23.03.2026
    \item Ημερομηνία Παράδοσης Αναφοράς: 30.03.2026
\end{itemize}
}

\newpage
\tableofcontents

\newpage
\section{Έλεγχος ορθής λειτουργίας της προσομοίωσης}

\subsection{Μεθοδολογία ελέγχου}

Πριν προχωρήσουμε στις κανονικές εκτελέσεις της προσομοίωσης για τις διάφορες τιμές του ρυθμού αφίξεων, είναι απαραίτητο να ελέγξουμε ότι ο βρόχος γεγονότων (\foreignlanguage{english}{event loop}) του προγράμματος λειτουργεί σωστά. Στο παρόν ερώτημα δεν μας ενδιαφέρει ακόμη η τελική ακρίβεια των εκτιμήσεων, αλλά η ορθότητα της δυναμικής του μοντέλου. Δηλαδή, πρέπει πρώτα να διαπιστώσουμε ότι το πρόγραμμα παράγει μόνο λογικές μεταβάσεις και ότι σέβεται πλήρως τους περιορισμούς της ουράς $\mathrm{M/M/1/10}$.

Για τον σκοπό αυτό ενεργοποιήσαμε ειδική λειτουργία \foreignlanguage{english}{debugging}, η οποία τυπώνει αναλυτικό \foreignlanguage{english}{trace} για τις πρώτες $30$ μεταβάσεις του συστήματος. Η επιλογή αυτή είναι απολύτως εύλογη, διότι μέσα σε ένα μικρό αλλά αντιπροσωπευτικό πλήθος βημάτων μπορούμε να ελέγξουμε εάν:

\begin{itemize}
    \item όταν το σύστημα είναι άδειο δεν εμφανίζεται ποτέ αναχώρηση,
    \item όταν το σύστημα είναι πλήρες οι νέες αφίξεις απορρίπτονται,
    \item η κατάσταση του συστήματος παραμένει πάντοτε στο επιτρεπτό διάστημα $0 \leq n \leq 10$,
    \item ο συνολικός αριθμός αφίξεων αυξάνεται μόνο όταν συμβαίνει προσπάθεια άφιξης,
    \item η νέα κατάσταση που παράγεται σε κάθε βήμα είναι συμβατή με το είδος του γεγονότος.
\end{itemize}

Στο \foreignlanguage{english}{debug} run χρησιμοποιήσαμε:
\[
\lambda = 5,\qquad \mu = 5,\qquad K = 10,\qquad T = 30,\qquad B = 0,
\]
με αρχική τιμή \foreignlanguage{english}{seed}
\[
\texttt{rand("seed", 1)}.
\]

Η επιλογή $\lambda = 5$ είναι ιδιαίτερα κατάλληλη για τον αρχικό αυτό έλεγχο, επειδή οδηγεί σε μία αρκετά ισορροπημένη συμπεριφορά του συστήματος. Έτσι, μέσα σε λίγες μόνο μεταβάσεις εμφανίζονται τόσο αφίξεις όσο και αναχωρήσεις, ενώ είναι αρκετά πιθανό να παρατηρήσουμε και γέμισμα της ουράς μέχρι τη μέγιστη χωρητικότητα, άρα και απορρίψεις πελατών. Με αυτόν τον τρόπο, ο έλεγχος ορθότητας δεν περιορίζεται σε μία ευνοϊκή περιοχή λειτουργίας του συστήματος, αλλά καλύπτει όλα τα κρίσιμα λογικά ενδεχόμενα.

\subsection{Υλοποίηση του ελέγχου στον κώδικα}

Η προσομοίωσή μας είναι \foreignlanguage{english}{event-driven}. Το σύστημα ξεκινά αρχικά άδειο, δηλαδή με:
\[
state = 0.
\]
Σε κάθε βήμα της προσομοίωσης παράγεται πρώτα ο χρόνος μέχρι το επόμενο γεγονός και στη συνέχεια αποφασίζεται αν το γεγονός αυτό είναι άφιξη ή αναχώρηση. Ειδικότερα:

\begin{itemize}
    \item αν το σύστημα βρίσκεται στην κατάσταση $state = 0$, το μόνο δυνατό επόμενο γεγονός είναι άφιξη,
    \item αν $state > 0$, τότε το επόμενο γεγονός είναι:
    \[
    \text{άφιξη με πιθανότητα } \frac{\lambda}{\lambda+\mu}
    \]
    και
    \[
    \text{αναχώρηση με πιθανότητα } \frac{\mu}{\lambda+\mu}.
    \]
\end{itemize}

Στη συνέχεια, η μετάβαση εκτελείται ως εξής:

\begin{enumerate}[label=\arabic*.]
    \item Αν συμβεί άφιξη και $state < K$, ο νέος πελάτης γίνεται δεκτός και η κατάσταση αυξάνεται κατά $1$.
    \item Αν συμβεί άφιξη και $state = K$, ο πελάτης απορρίπτεται και η κατάσταση παραμένει αμετάβλητη.
    \item Αν συμβεί αναχώρηση, αφαιρείται ο παλαιότερος πελάτης από το σύστημα και η κατάσταση μειώνεται κατά $1$.
\end{enumerate}

Παράλληλα, για τα πρώτα $30$ βήματα τυπώνονται στην οθόνη τα ακόλουθα στοιχεία:

\begin{itemize}
    \item ο αριθμός του βήματος,
    \item η τρέχουσα κατάσταση πριν από το γεγονός,
    \item το είδος του γεγονότος,
    \item ο συνολικός αριθμός αφίξεων που έχουν συμβεί μέχρι εκείνη τη στιγμή,
    \item η νέα κατάσταση μετά τη μετάβαση.
\end{itemize}

Αξίζει να σημειωθεί ότι στον συνολικό αριθμό αφίξεων προσμετρώνται όλες οι προσπάθειες άφιξης, είτε αυτές γίνονται δεκτές είτε απορρίπτονται λόγω πλήρους συστήματος. Αυτό είναι το σωστό, διότι η πιθανότητα απόρριψης ορίζεται ως λόγος απορριφθέντων πελατών προς το σύνολο των προσπαθειών άφιξης.

\subsection{Έξοδος \foreignlanguage{english}{debugging} για τις πρώτες 30 μεταβάσεις}

Παρακάτω παρατίθεται αυτούσια η έξοδος του προγράμματος για το \foreignlanguage{english}{debug} run:

{
\begin{verbatim}
DEBUG TRACE - first 30 transitions
lambda = 5, T = 30, B = 0, seed = 1
============================================================
Step  1 | state= 0 | event=arrival (accepted) | arrivals= 1 | new state= 1
Step  2 | state= 1 | event=departure        | arrivals= 1 | new state= 0
Step  3 | state= 0 | event=arrival (accepted) | arrivals= 2 | new state= 1
Step  4 | state= 1 | event=departure        | arrivals= 2 | new state= 0
Step  5 | state= 0 | event=arrival (accepted) | arrivals= 3 | new state= 1
Step  6 | state= 1 | event=arrival (accepted) | arrivals= 4 | new state= 2
Step  7 | state= 2 | event=arrival (accepted) | arrivals= 5 | new state= 3
Step  8 | state= 3 | event=arrival (accepted) | arrivals= 6 | new state= 4
Step  9 | state= 4 | event=arrival (accepted) | arrivals= 7 | new state= 5
Step 10 | state= 5 | event=arrival (accepted) | arrivals= 8 | new state= 6
Step 11 | state= 6 | event=arrival (accepted) | arrivals= 9 | new state= 7
Step 12 | state= 7 | event=departure        | arrivals= 9 | new state= 6
Step 13 | state= 6 | event=arrival (accepted) | arrivals=10 | new state= 7
Step 14 | state= 7 | event=arrival (accepted) | arrivals=11 | new state= 8
Step 15 | state= 8 | event=arrival (accepted) | arrivals=12 | new state= 9
Step 16 | state= 9 | event=arrival (accepted) | arrivals=13 | new state=10
Step 17 | state=10 | event=arrival (rejected) | arrivals=14 | new state=10
Step 18 | state=10 | event=arrival (rejected) | arrivals=15 | new state=10
Step 19 | state=10 | event=arrival (rejected) | arrivals=16 | new state=10
Step 20 | state=10 | event=departure        | arrivals=16 | new state= 9
Step 21 | state= 9 | event=arrival (accepted) | arrivals=17 | new state=10
Step 22 | state=10 | event=arrival (rejected) | arrivals=18 | new state=10
Step 23 | state=10 | event=departure        | arrivals=18 | new state= 9
Step 24 | state= 9 | event=departure        | arrivals=18 | new state= 8
Step 25 | state= 8 | event=departure        | arrivals=18 | new state= 7
Step 26 | state= 7 | event=arrival (accepted) | arrivals=19 | new state= 8
Step 27 | state= 8 | event=arrival (accepted) | arrivals=20 | new state= 9
Step 28 | state= 9 | event=arrival (accepted) | arrivals=21 | new state=10
Step 29 | state=10 | event=departure        | arrivals=21 | new state= 9
Step 30 | state= 9 | event=departure        | arrivals=21 | new state= 8
\end{verbatim}
}

\subsection{Σχολιασμός της εξόδου και επαλήθευση ορθότητας}

Από το παραπάνω \foreignlanguage{english}{trace} προκύπτουν σαφώς αρκετές ενδείξεις ότι η προσομοίωση λειτουργεί σωστά.

Αρχικά, παρατηρούμε ότι στο πρώτο βήμα, όπου το σύστημα είναι άδειο ($state = 0$), το επόμενο γεγονός είναι υποχρεωτικά άφιξη:
\[
0 \to 1.
\]
Αυτό είναι απολύτως σωστό, διότι σε άδειο σύστημα δεν μπορεί να συμβεί αναχώρηση. Η ίδια λογική επιβεβαιώνεται και στα πρώτα βήματα του \foreignlanguage{english}{trace}, όπου κάθε φορά που το σύστημα επιστρέφει στην κατάσταση $0$, η επόμενη μετάβαση είναι και πάλι άφιξη.

Στη συνέχεια, παρατηρούμε ότι το σύστημα γεμίζει σταδιακά μέσω διαδοχικών αποδεκτών αφίξεων. Από το βήμα $6$ μέχρι και το βήμα $16$, η κατάσταση αυξάνεται από $1$ μέχρι $10$, δηλαδή:
\[
1 \to 2 \to 3 \to \cdots \to 10.
\]
Η εξέλιξη αυτή είναι πλήρως συμβατή με τη λογική του μοντέλου, αφού κάθε αποδεκτή άφιξη αυξάνει τον αριθμό πελατών στο σύστημα κατά μία μονάδα.

Το σημαντικότερο ίσως σημείο ελέγχου εμφανίζεται στα βήματα $17$, $18$ και $19$. Εκεί βλέπουμε ότι, ενώ το σύστημα βρίσκεται ήδη στην κατάσταση $10$, συμβαίνουν νέες προσπάθειες άφιξης, οι οποίες καταγράφονται κανονικά ως αφίξεις, αλλά απορρίπτονται:
\[
10 \to 10.
\]
Αυτό ακριβώς πρέπει να συμβαίνει σε ουρά $\mathrm{M/M/1/10}$, αφού η χωρητικότητα του συστήματος είναι πεπερασμένη και ίση με $10$. Επομένως, ο μηχανισμός \foreignlanguage{english}{blocking} έχει υλοποιηθεί σωστά.

Επιπλέον, παρατηρούμε ότι ο μετρητής \texttt{arrivals} αυξάνεται μόνο όταν συμβαίνει προσπάθεια άφιξης και παραμένει σταθερός όταν συμβαίνει αναχώρηση. Για παράδειγμα:

\begin{itemize}
    \item στο βήμα $16$ έχουμε \texttt{arrivals = 13},
    \item στα βήματα $17$--$19$ έχουμε τρεις απορριπτόμενες αφίξεις, άρα ο μετρητής αυξάνεται μέχρι \texttt{arrivals = 16},
    \item στο βήμα $20$ συμβαίνει αναχώρηση και ο μετρητής παραμένει \texttt{16}.
\end{itemize}

Το παραπάνω είναι επίσης απολύτως ορθό, αφού η απόρριψη ενός πελάτη δεν αναιρεί το γεγονός ότι πραγματοποιήθηκε προσπάθεια άφιξης.

Παρατηρούμε ακόμη ότι οι αναχωρήσεις μειώνουν πάντοτε την κατάσταση του συστήματος κατά μία μονάδα:
\[
10 \to 9,\qquad 9 \to 8,\qquad 8 \to 7,
\]
ενώ δεν εμφανίζεται ποτέ μετάβαση που να οδηγεί σε αρνητική κατάσταση. Αντίστοιχα, δεν εμφανίζεται ποτέ αποδεκτή άφιξη όταν το σύστημα είναι ήδη πλήρες. Συνεπώς, η προσομοίωση διατηρεί την κατάσταση αυστηρά μέσα στο επιτρεπτό διάστημα:
\[
0 \leq state \leq 10.
\]

Ένα ακόμη πολύ σημαντικό στοιχείο είναι ότι στο συγκεκριμένο \foreignlanguage{english}{debug} run εμφανίζονται και οι τρεις βασικές κατηγορίες γεγονότων που θέλαμε να ελέγξουμε:

\begin{enumerate}[label=\roman*)]
    \item αποδεκτές αφίξεις,
    \item αναχωρήσεις,
    \item απορρίψεις λόγω πλήρους συστήματος.
\end{enumerate}

Άρα, μέσα στις πρώτες $30$ μεταβάσεις ελέγξαμε εμπειρικά όλα τα κρίσιμα λογικά μονοπάτια του \foreignlanguage{english}{event loop}. Αυτό είναι ιδιαίτερα χρήσιμο, διότι δεν επαληθεύεται απλώς μία τοπική περιοχή λειτουργίας του κώδικα, αλλά επιβεβαιώνεται η ορθότητα των κανόνων μετάβασης τόσο σε χαμηλό όσο και σε υψηλό φόρτο του συστήματος.

\subsection{Σχόλιο για τη μεθοδολογία της υλοποίησης}

Ο αρχικός βοηθητικός κώδικας που μας δόθηκε διατηρούσε τη βασική λογική των μεταβάσεων από κατάσταση σε κατάσταση, δηλαδή το ποια γεγονότα επιτρέπονται και με ποιον τρόπο ενημερώνεται το \texttt{state}. Στη δική μας υλοποίηση διατηρήσαμε αυτή τη βασική λογική, αλλά οργανώσαμε τον προσομοιωτή σε πλήρη \foreignlanguage{english}{event-driven} μορφή, ώστε να υποστηρίζει όχι μόνο τον έλεγχο του βρόχου γεγονότων, αλλά και τον υπολογισμό όλων των ζητούμενων μεγεθών της άσκησης.

Συγκεκριμένα, εκτός από τον έλεγχο των μεταβάσεων, ο κώδικας:

\begin{itemize}
    \item παράγει τον χρόνο μέχρι το επόμενο γεγονός,
    \item διατηρεί τις χρονικές στιγμές άφιξης των πελατών,
    \item υπολογίζει \foreignlanguage{english}{time-average} εργοδικές πιθανότητες,
    \item υπολογίζει τις εργοδικές πιθανότητες τόσο με όλες τις μεταβάσεις ($P$) όσο και μετά την αφαίρεση του transient ($P_B$),
    \item υπολογίζει τον μέσο αριθμό πελατών $L$,
    \item υπολογίζει τον μέσο χρόνο παραμονής $W$,
    \item υπολογίζει τα ίδια μεγέθη τόσο με όλες τις μεταβάσεις όσο και μετά την αφαίρεση της αρχικής μεταβατικής περιόδου,
    \item αναπαριστά πλέον στο ίδιο διάγραμμα τις δύο εκδοχές των εργοδικών πιθανοτήτων, ώστε να είναι άμεση η σύγκρισή τους.
\end{itemize}

Ωστόσο, στο παρόν πρώτο ερώτημα ο έλεγχος ορθότητας επικεντρώνεται κυρίως στη συμπεριφορά του \foreignlanguage{english}{event loop}. Αν το \foreignlanguage{english}{trace} περιείχε παράλογες μεταβάσεις, όπως αναχώρηση από άδειο σύστημα, αποδεκτή άφιξη σε πλήρες σύστημα ή κατάσταση εκτός του διαστήματος $[0,10]$, τότε ολόκληρη η προσομοίωση θα έπρεπε να θεωρηθεί λανθασμένη και δεν θα είχε νόημα να προχωρήσουμε στην παραγωγή αριθμητικών αποτελεσμάτων και διαγραμμάτων.

\newpage
\section{Εκτέλεση της προσομοίωσης για διαφορετικές τιμές του ρυθμού άφιξης}

\subsection{Διάταξη των εκτελέσεων και σημεία της υλοποίησης}

Αφού ολοκληρώθηκε επιτυχώς ο έλεγχος ορθής λειτουργίας του βρόχου γεγονότων, απενεργοποιήσαμε το \foreignlanguage{english}{debug print} και προχωρήσαμε στις κανονικές εκτελέσεις της προσομοίωσης. Σύμφωνα με την εκφώνηση, θεωρήσαμε σταθερό ρυθμό εξυπηρέτησης
\[
\mu = 5 \ \text{πελάτες/min},
\]
μέγιστη χωρητικότητα
\[
K = 10,
\]
και τρέξαμε την προσομοίωση για τις τρεις τιμές του ρυθμού αφίξεων
\[
\lambda \in \{1,5,10\},
\]
ώστε να καλυφθούν οι περιπτώσεις χαμηλού, οριακού και υψηλού φόρτου λειτουργίας του συστήματος.

Επιπλέον, για κάθε τιμή του $\lambda$ χρησιμοποιήσαμε και τα τρία ζητούμενα ζεύγη τιμών:
\[
(T,B) = (100000,5000),\qquad (1000000,20000),\qquad (10000000,120000).
\]
Κατά συνέπεια, εκτελέστηκαν συνολικά $9$ κανονικά \foreignlanguage{english}{simulation runs}, δηλαδή τρία \foreignlanguage{english}{runs} για κάθε τιμή του $\lambda$.

Σημαντικό σημείο της υλοποίησης είναι ότι σε κάθε εκτέλεση επαναφέραμε την ίδια αρχική τιμή της γεννήτριας τυχαίων αριθμών, δηλαδή το ίδιο \foreignlanguage{english}{seed}. Με τον τρόπο αυτό εξασφαλίζεται ότι η σύγκριση μεταξύ διαφορετικών τιμών του $\lambda$ και διαφορετικών ζευγών $(T,B)$ γίνεται με συνεπή τρόπο, χωρίς να εισάγονται επιπλέον αποκλίσεις λόγω διαφορετικών ακολουθιών τυχαίων αριθμών.

Σε επίπεδο κώδικα, η εκτέλεση των πειραμάτων οργανώθηκε μέσω δύο εμφωλευμένων βρόχων:
\begin{itemize}
    \item ο εξωτερικός βρόχος διατρέχει τα τρία ζεύγη $(T,B)$,
    \item ο εσωτερικός βρόχος διατρέχει τις τρεις τιμές του $\lambda$.
\end{itemize}

Για κάθε συνδυασμό τιμών, η συνάρτηση της προσομοίωσης επιστρέφει μία δομή αποτελεσμάτων που περιέχει:
\begin{itemize}
    \item τις εργοδικές πιθανότητες των καταστάσεων με όλες τις μεταβάσεις,
    \item τις εργοδικές πιθανότητες μετά την αφαίρεση των πρώτων $B$ μεταβάσεων,
    \item τον μέσο αριθμό πελατών στο σύστημα $L$,
    \item την πιθανότητα απόρριψης $P_{rej}$,
    \item τον μέσο χρόνο παραμονής $W$,
    \item τις αντίστοιχες μετα-μεταβατικές εκτιμήσεις $L_B$, $P_{rej,B}$ και $W_B$,
    \item τον χρόνο εκτέλεσης της προσομοίωσης.
\end{itemize}

Παράλληλα, για να είναι δυνατή η σχεδίαση των ζητούμενων γραφημάτων, αποθηκεύεται κατά τη διάρκεια κάθε εκτέλεσης η εξέλιξη της εκτίμησης του μέσου πλήθους πελατών $L$ ως προς τον αριθμό των μεταβάσεων. Έτσι, στο τέλος κάθε \foreignlanguage{english}{run} παράγονται δύο κατηγορίες γραφημάτων:
\begin{enumerate}[label=\arabic*.]
    \item διαγράμματα με τις εργοδικές πιθανότητες των καταστάσεων με όλες τις μεταβάσεις και μετά την αφαίρεση του transient στο ίδιο γράφημα,
    \item διαγράμματα με την εξέλιξη της εκτίμησης του $L$ με όλα τα δεδομένα και μετά την αφαίρεση των πρώτων $B$ μεταβάσεων.
\end{enumerate}

Στην παρούσα ενότητα παραθέτουμε τα ακατέργαστα αριθμητικά αποτελέσματα και τα αντίστοιχα διαγράμματα. Η ουσιαστική ερμηνεία και ο σχολιασμός τους θα γίνει στην επόμενη ενότητα.

\subsection{Έξοδος προγράμματος για $(T,B) = (100000,5000)$}

\subsubsection*{Περίπτωση $\lambda = 1$}

{
\begin{verbatim}
============================================================
Running simulation for lambda = 1, T = 100000, B = 5000
============================================================
lambda = 1, mu = 5, K = 10, T = 100000, B = 5000
Runtime = 6.6001 sec

--- Using all transitions ---
Mean number in system L       = 0.2511296521
Rejection probability P_rej   = 0.0000000000
Mean time in system W         = 0.2505772043
State probabilities P         = [ 0.79945865 0.16008330 0.03232467 0.00655883
                                  0.00122766 0.00027151 0.00007539 0.00000000
                                  0.00000000 0.00000000 0.00000000 ]

--- Ignoring first B transitions ---
Mean number in system L_B     = 0.2509471723
Rejection probability P_rej_B = 0.0000000000
Mean time in system W_B       = 0.2506592621
State probabilities P_B       = [ 0.79959641 0.15997915 0.03231434 0.00654240
                                  0.00120567 0.00028276 0.00007928 0.00000000
                                  0.00000000 0.00000000 0.00000000 ]

Relative difference in L      = 0.072664 %
Relative difference in P_rej  = 0.000000 %
Relative difference in W      = 0.032748 %
\end{verbatim}
}

\subsubsection*{Περίπτωση $\lambda = 5$}

{
\begin{verbatim}
============================================================
Running simulation for lambda = 5, T = 100000, B = 5000
============================================================
lambda = 5, mu = 5, K = 10, T = 100000, B = 5000
Runtime = 6.9150 sec

--- Using all transitions ---
Mean number in system L       = 5.0669430209
Rejection probability P_rej   = 0.0955210238
Mean time in system W         = 1.1218601440
State probabilities P         = [ 0.08966699 0.08890558 0.08922529 0.08838208
                                  0.08898398 0.09082476 0.09041571 0.09071322
                                  0.09425913 0.09341150 0.09521176 ]

--- Ignoring first B transitions ---
Mean number in system L_B     = 5.1071053828
Rejection probability P_rej_B = 0.0968930911
Mean time in system W_B       = 1.1319838096
State probabilities P_B       = [ 0.08845269 0.08741447 0.08785467 0.08707792
                                  0.08838430 0.09068520 0.09099106 0.09164490
                                  0.09596627 0.09469143 0.09683710 ]

Relative difference in L      = 0.792635 %
Relative difference in P_rej  = 1.436403 %
Relative difference in W      = 0.902400 %
\end{verbatim}
}

\subsubsection*{Περίπτωση $\lambda = 10$}

{
\begin{verbatim}
============================================================
Running simulation for lambda = 10, T = 100000, B = 5000
============================================================
lambda = 10, mu = 5, K = 10, T = 100000, B = 5000
Runtime = 5.8064 sec

--- Using all transitions ---
Mean number in system L       = 9.0373462585
Rejection probability P_rej   = 0.5063846947
Mean time in system W         = 1.8266043984
State probabilities P         = [ 0.00081372 0.00144580 0.00216852 0.00373982
                                  0.00677244 0.01378517 0.02777216 0.06007758
                                  0.12141213 0.25427131 0.50774136 ]

--- Ignoring first B transitions ---
Mean number in system L_B     = 9.0350260323
Rejection probability P_rej_B = 0.5063205560
Mean time in system W_B       = 1.8278095613
State probabilities P_B       = [ 0.00082343 0.00149152 0.00226441 0.00376726
                                  0.00672410 0.01396358 0.02808737 0.05998343
                                  0.12131367 0.25374031 0.50784093 ]

Relative difference in L      = 0.025674 %
Relative difference in P_rej  = 0.012666 %
Relative difference in W      = 0.065978 %
\end{verbatim}
}

\subsection{Έξοδος προγράμματος για $(T,B) = (1000000,20000)$}

\subsubsection*{Περίπτωση $\lambda = 1$}

{
\begin{verbatim}
============================================================
Running simulation for lambda = 1, T = 1000000, B = 20000
============================================================
lambda = 1, mu = 5, K = 10, T = 1000000, B = 20000
Runtime = 65.9007 sec

--- Using all transitions ---
Mean number in system L       = 0.2505027289
Rejection probability P_rej   = 0.0000000000
Mean time in system W         = 0.2503982626
State probabilities P         = [ 0.79980439 0.15989668 0.03226471 0.00646411
                                  0.00125921 0.00023789 0.00005671 0.00001337
                                  0.00000214 0.00000078 0.00000000 ]

--- Ignoring first B transitions ---
Mean number in system L_B     = 0.2503824648
Rejection probability P_rej_B = 0.0000000000
Mean time in system W_B       = 0.2504195813
State probabilities P_B       = [ 0.79990402 0.15981203 0.03225954 0.00645092
                                  0.00126277 0.00023721 0.00005691 0.00001364
                                  0.00000218 0.00000079 0.00000000 ]

Relative difference in L      = 0.048009 %
Relative difference in P_rej  = 0.000000 %
Relative difference in W      = 0.008514 %
\end{verbatim}
}

\subsubsection*{Περίπτωση $\lambda = 5$}

{
\begin{verbatim}
============================================================
Running simulation for lambda = 5, T = 1000000, B = 20000
============================================================
lambda = 5, mu = 5, K = 10, T = 1000000, B = 20000
Runtime = 67.4362 sec

--- Using all transitions ---
Mean number in system L       = 5.0026245648
Rejection probability P_rej   = 0.0898944867
Mean time in system W         = 1.1012411221
State probabilities P         = [ 0.08999344 0.09049205 0.09038443 0.09076960
                                  0.09155982 0.09187900 0.09182116 0.09179528
                                  0.09157453 0.08997642 0.08975426 ]

--- Ignoring first B transitions ---
Mean number in system L_B     = 5.0080060496
Rejection probability P_rej_B = 0.0901637587
Mean time in system W_B       = 1.1026525045
State probabilities P_B       = [ 0.08983444 0.09029272 0.09023266 0.09062306
                                  0.09149535 0.09182429 0.09177938 0.09192607
                                  0.09181162 0.09017976 0.09000064 ]

Relative difference in L      = 0.107573 %
Relative difference in P_rej  = 0.299542 %
Relative difference in W      = 0.128163 %
\end{verbatim}
}

\subsubsection*{Περίπτωση $\lambda = 10$}

{
\begin{verbatim}
============================================================
Running simulation for lambda = 10, T = 1000000, B = 20000
============================================================
lambda = 10, mu = 5, K = 10, T = 1000000, B = 20000
Runtime = 59.6151 sec

--- Using all transitions ---
Mean number in system L       = 9.0193533021
Rejection probability P_rej   = 0.5019511910
Mean time in system W         = 1.8090288162
State probabilities P         = [ 0.00047210 0.00093474 0.00183579 0.00375632
                                  0.00744585 0.01515244 0.03045462 0.06143927
                                  0.12444488 0.25106915 0.50299484 ]

--- Ignoring first B transitions ---
Mean number in system L_B     = 9.0198466095
Rejection probability P_rej_B = 0.5021108010
Mean time in system W_B       = 1.8096164018
State probabilities P_B       = [ 0.00046298 0.00089190 0.00181307 0.00376460
                                  0.00744556 0.01519379 0.03049994 0.06138163
                                  0.12443250 0.25098774 0.50312628 ]

Relative difference in L      = 0.005469 %
Relative difference in P_rej  = 0.031798 %
Relative difference in W      = 0.032481 %
\end{verbatim}
}

\subsection{Έξοδος προγράμματος για $(T,B) = (10000000,120000)$}

\subsubsection*{Περίπτωση $\lambda = 1$}

{
\begin{verbatim}
============================================================
Running simulation for lambda = 1, T = 10000000, B = 120000
============================================================
lambda = 1, mu = 5, K = 10, T = 10000000, B = 120000
Runtime = 635.6335 sec

--- Using all transitions ---
Mean number in system L       = 0.2504273515
Rejection probability P_rej   = 0.0000000000
Mean time in system W         = 0.2502386518
State probabilities P         = [ 0.79975162 0.16010075 0.03211215 0.00643726
                                  0.00127917 0.00025538 0.00005141 0.00001010
                                  0.00000168 0.00000048 0.00000000 ]

--- Ignoring first B transitions ---
Mean number in system L_B     = 0.2504118097
Rejection probability P_rej_B = 0.0000000000
Mean time in system W_B       = 0.2502303841
State probabilities P_B       = [ 0.79976174 0.16009561 0.03210828 0.00643548
                                  0.00127993 0.00025533 0.00005124 0.00001020
                                  0.00000170 0.00000048 0.00000000 ]

Relative difference in L      = 0.006206 %
Relative difference in P_rej  = 0.000000 %
Relative difference in W      = 0.003304 %
\end{verbatim}
}

\subsubsection*{Περίπτωση $\lambda = 5$}

{
\begin{verbatim}
============================================================
Running simulation for lambda = 5, T = 10000000, B = 120000
============================================================
lambda = 5, mu = 5, K = 10, T = 10000000, B = 120000
Runtime = 639.4556 sec

--- Using all transitions ---
Mean number in system L       = 5.0076954577
Rejection probability P_rej   = 0.0911728040
Mean time in system W         = 1.1017529664
State probabilities P         = [ 0.09040039 0.09059619 0.09068113 0.09101954
                                  0.09087455 0.09094612 0.09094979 0.09098016
                                  0.09122787 0.09117577 0.09114848 ]

--- Ignoring first B transitions ---
Mean number in system L_B     = 5.0075222213
Rejection probability P_rej_B = 0.0911554245
Mean time in system W_B       = 1.1016549734
State probabilities P_B       = [ 0.09039147 0.09060094 0.09068645 0.09103910
                                  0.09088581 0.09093975 0.09095386 0.09098918
                                  0.09121022 0.09117228 0.09113094 ]

Relative difference in L      = 0.003459 %
Relative difference in P_rej  = 0.019062 %
Relative difference in W      = 0.008894 %
\end{verbatim}
}

\subsubsection*{Περίπτωση $\lambda = 10$}

{
\begin{verbatim}
============================================================
Running simulation for lambda = 10, T = 10000000, B = 120000
============================================================
lambda = 10, mu = 5, K = 10, T = 10000000, B = 120000
Runtime = 626.7589 sec

--- Using all transitions ---
Mean number in system L       = 9.0065989686
Rejection probability P_rej   = 0.5004547366
Mean time in system W         = 1.8020232187
State probabilities P         = [ 0.00046642 0.00096135 0.00192156 0.00390119
                                  0.00784720 0.01562832 0.03114148 0.06239525
                                  0.12498207 0.25046324 0.50029192 ]

--- Ignoring first B transitions ---
Mean number in system L_B     = 9.0062655890
Rejection probability P_rej_B = 0.5003975920
Mean time in system W_B       = 1.8017936156
State probabilities P_B       = [ 0.00046348 0.00095733 0.00192049 0.00390357
                                  0.00785770 0.01564922 0.03117419 0.06242214
                                  0.12502019 0.25039893 0.50023277 ]

Relative difference in L      = 0.003702 %
Relative difference in P_rej  = 0.011419 %
Relative difference in W      = 0.012741 %
\end{verbatim}
}

\subsection{Παραγόμενα διαγράμματα}

Στη συνέχεια παραθέτουμε τα έξι διαγράμματα που παρήχθησαν από τις κανονικές εκτελέσεις της προσομοίωσης. Για κάθε ζεύγος $(T,B)$ παρουσιάζονται:
\begin{itemize}
    \item το διάγραμμα των εργοδικών πιθανοτήτων με όλες τις μεταβάσεις και μετά την αφαίρεση του transient στο ίδιο γράφημα,
    \item το διάγραμμα εξέλιξης της εκτίμησης του μέσου αριθμού πελατών $L$ ως προς τον αριθμό των μεταβάσεων.
\end{itemize}

\begin{figure}[p]
    \centering
    \includegraphics[width=\textwidth]{t_100000_b_5000_prob.png}
    \caption{Εργοδικές πιθανότητες με όλες τις μεταβάσεις και μετά την αφαίρεση του transient για $(T,B)=(100000,5000)$.}
\end{figure}

\begin{figure}[p]
    \centering
    \includegraphics[width=\textwidth]{t_100000_b_5000_diag.png}
    \caption{Εξέλιξη της εκτίμησης του μέσου πλήθους πελατών $L$ για $(T,B)=(100000,5000)$.}
\end{figure}

\begin{figure}[p]
    \centering
    \includegraphics[width=\textwidth]{t_1000000_b_20000_prob.png}
    \caption{Εργοδικές πιθανότητες με όλες τις μεταβάσεις και μετά την αφαίρεση του transient για $(T,B)=(1000000,20000)$.}
\end{figure}

\begin{figure}[p]
    \centering
    \includegraphics[width=\textwidth]{t_1000000_b_20000_diag.png}
    \caption{Εξέλιξη της εκτίμησης του μέσου πλήθους πελατών $L$ για $(T,B)=(1000000,20000)$.}
\end{figure}

\begin{figure}[p]
    \centering
    \includegraphics[width=\textwidth]{t_10000000_b_120000_prob.png}
    \caption{Εργοδικές πιθανότητες με όλες τις μεταβάσεις και μετά την αφαίρεση του transient για $(T,B)=(10000000,120000)$.}
\end{figure}

\begin{figure}[p]
    \centering
    \includegraphics[width=\textwidth]{t_10000000_b_120000_diag.png}
    \caption{Εξέλιξη της εκτίμησης του μέσου πλήθους πελατών $L$ για $(T,B)=(10000000,120000)$.}
\end{figure}

\newpage
\section{Σχολιασμός των αποτελεσμάτων}

\subsection{Θεωρητικό πλαίσιο και αναμενόμενη συμπεριφορά}

Πριν σχολιάσουμε αναλυτικά τα αριθμητικά αποτελέσματα της προσομοίωσης, είναι χρήσιμο να διατυπώσουμε συνοπτικά τι αναμένουμε θεωρητικά από ένα σύστημα $\mathrm{M/M/1/10}$. Πρόκειται για μία ουρά με αφίξεις Poisson, εκθετικούς χρόνους εξυπηρέτησης, έναν εξυπηρετητή και συνολική χωρητικότητα $K=10$. Επομένως, το σύστημα μπορεί να βρίσκεται μόνο στις καταστάσεις
\[
n=0,1,2,\ldots,10,
\]
ενώ όταν το σύστημα γεμίσει πλήρως, δηλαδή όταν $n=10$, κάθε νέα άφιξη απορρίπτεται.

Αν θέσουμε
\[
\rho=\frac{\lambda}{\mu},
\]
τότε για το πεπερασμένο σύστημα $\mathrm{M/M/1/K}$ με $K=10$, οι εργοδικές πιθανότητες των καταστάσεων δίνονται θεωρητικά από:
\[
P_n=\frac{(1-\rho)\rho^n}{1-\rho^{K+1}}, \qquad n=0,1,\ldots,K,\qquad \rho\neq 1.
\]
Στην ειδική περίπτωση $\rho=1$, δηλαδή όταν $\lambda=\mu$, όλες οι καταστάσεις είναι ισοπίθανες:
\[
P_n=\frac{1}{K+1}=\frac{1}{11}, \qquad n=0,1,\ldots,10.
\]

Επιπλέον, σε ένα σύστημα πεπερασμένης χωρητικότητας η πιθανότητα απόρριψης ενός πελάτη είναι ίση με την εργοδική πιθανότητα της πλήρους κατάστασης:
\[
P_{rej}=P_{10}.
\]
Αν ο ενεργός ρυθμός αφίξεων στο σύστημα είναι
\[
\lambda_{\text{eff}}=\lambda(1-P_{rej}),
\]
τότε, από τον νόμο του Little, ο μέσος χρόνος παραμονής στο σύστημα προκύπτει από:
\[
W=\frac{L}{\lambda_{\text{eff}}},
\]
όπου $L$ είναι ο μέσος αριθμός πελατών στο σύστημα.

Με βάση τα παραπάνω, αναμένουμε τα εξής:

\begin{itemize}
    \item για $\lambda=1$ και $\mu=5$, δηλαδή $\rho=0.2$, το σύστημα θα είναι σαφώς \foreignlanguage{english}{underloaded}, άρα θα παραμένει τις περισσότερες φορές σχεδόν άδειο και η πιθανότητα απόρριψης θα είναι πρακτικά μηδενική,
    \item για $\lambda=5$ και $\mu=5$, δηλαδή $\rho=1$, το σύστημα θα βρίσκεται σε οριακή κατάσταση φόρτου, οπότε αναμένουμε αρκετά πιο απλωμένη ή και ομοιόμορφη κατανομή στις καταστάσεις και αισθητή πιθανότητα blocking,
    \item για $\lambda=10$ και $\mu=5$, δηλαδή $\rho=2$, το σύστημα θα βρίσκεται σε έντονα φορτισμένη περιοχή λειτουργίας, με συχνή προσέγγιση ή παραμονή στη μέγιστη χωρητικότητα, μεγάλο μέσο αριθμό πελατών και υψηλή πιθανότητα απόρριψης.
\end{itemize}

Όπως θα φανεί αμέσως παρακάτω, τα αποτελέσματα της προσομοίωσης συμφωνούν σε πολύ μεγάλο βαθμό με αυτή τη θεωρητική εικόνα.

\subsection{Συνοπτική εικόνα των αποτελεσμάτων}

Για να έχουμε μία συνολική και συγκεντρωτική εικόνα, παραθέτουμε πρώτα τα αποτελέσματα της πιο εκτενούς εκτέλεσης, δηλαδή για
\[
(T,B)=(10^7,120000),
\]
καθώς εκεί οι σχετικές διαφορές μεταξύ με transient και χωρίς transient είναι οι μικρότερες και οι εκτιμήσεις έχουν σταθεροποιηθεί περισσότερο.

\begin{table}[h!]
\centering
\caption{Συγκεντρωτικά αποτελέσματα για $(T,B)=(10^7,120000)$.}
\begin{tabular}{|c|P{4.7cm}|P{4.7cm}|P{4.7cm}|}
\hline
$\lambda$ & $(L,\;L_B)$ & $(P_{rej},\;P_{rej,B})$ & $(W,\;W_B)$ \\ \hline
$1$  & $(0.2504273515,\;0.2504118097)$ & $(0.0000000000,\;0.0000000000)$ & $(0.2502386518,\;0.2502303841)$ \\ \hline
$5$  & $(5.0076954577,\;5.0075222213)$ & $(0.0911728040,\;0.0911554245)$ & $(1.1017529664,\;1.1016549734)$ \\ \hline
$10$ & $(9.0065989686,\;9.0062655890)$ & $(0.5004547366,\;0.5003975920)$ & $(1.8020232187,\;1.8017936156)$ \\ \hline
\end{tabular}
\end{table}

Ήδη από τον παραπάνω πίνακα φαίνεται ξεκάθαρα η βασική τάση: καθώς ο ρυθμός άφιξης $\lambda$ αυξάνεται, αυξάνεται και ο μέσος αριθμός πελατών στο σύστημα, αυξάνεται η πιθανότητα απόρριψης και αυξάνεται ο μέσος χρόνος παραμονής των πελατών που τελικά γίνονται δεκτοί στο σύστημα. Ταυτόχρονα, παρατηρούμε ότι στις μεγάλες εκτελέσεις οι εκτιμήσεις με όλες τις μεταβάσεις και μετά την αφαίρεση του transient είναι πρακτικά ταυτόσημες.

\subsection{Άμεση σύγκριση των εργοδικών πιθανοτήτων με και χωρίς transient}

Καθώς στο ίδιο εκάστοτε διάγραμμα της προηγούμενης ενότητας αναπαρίστανται οι εργοδικές πιθανότητες και πριν και μετά την αφαίρεση του transient, αυτό μας επιτρέπει να συγκρίνουμε οπτικά πόσο επηρεάζει η αρχική μεταβατική περίοδος την τελική εκτίμηση της κατανομής των καταστάσεων.

Από τα grouped bar diagrams φαίνεται ότι:

\begin{itemize}
    \item για $\lambda=1$ οι μπάρες των $P$ και $P_B$ είναι ήδη σχεδόν ταυτόσημες ακόμη και για $(T,B)=(100000,5000)$,
    \item για $\lambda=5$ η διαφορά είναι περισσότερο ορατή στο μικρότερο run, ιδίως στις υψηλότερες καταστάσεις, όμως μειώνεται αισθητά για $T=10^6$ και γίνεται ουσιαστικά αμελητέα για $T=10^7$,
    \item για $\lambda=10$ οι δύο κατανομές είναι επίσης πολύ κοντά μεταξύ τους, με μικρές μόνο αποκλίσεις που μειώνονται περαιτέρω όσο αυξάνει το $T$.
\end{itemize}

Η οπτική αυτή σύγκριση επιβεβαιώνει άμεσα ό,τι έδειχναν ήδη και οι αριθμητικές σχετικές διαφορές: η αφαίρεση του transient έχει ορατότερη επίδραση κυρίως στην οριακή περίπτωση $\lambda=5$, ενώ στις άλλες δύο περιπτώσεις η διαφορά είναι εξαρχής πολύ μικρότερη.

\subsection{Επίδραση της αφαίρεσης της αρχικής μεταβατικής περιόδου}

Η πρώτη ερώτηση που καλούμαστε να απαντήσουμε είναι κατά πόσο η αφαίρεση των πρώτων $B$ μεταβάσεων επηρεάζει τις εκτιμήσεις των μεγεθών. Αυτός ακριβώς είναι ο σκοπός της έννοιας \foreignlanguage{english}{transient period}: το σύστημα ξεκινά από την τεχνητή αρχική κατάσταση $n(0)=0$, η οποία δεν αντιστοιχεί γενικά στη μόνιμη λειτουργία του συστήματος. Άρα, αν χρησιμοποιήσουμε και τα πρώτα βήματα στην εκτίμηση των μετρικών, εισάγουμε μία αρχική μεροληψία.

Στα αποτελέσματά μας παρατηρούμε ότι η αφαίρεση του transient επηρεάζει με διαφορετικό τρόπο τις τρεις τιμές του $\lambda$.

\subsubsection*{Περίπτωση $\lambda=1$}

Για $\lambda=1$ οι διαφορές μεταξύ των δύο τρόπων υπολογισμού είναι πολύ μικρές σε όλα τα ζεύγη $(T,B)$. Συγκεκριμένα, για:
\begin{itemize}
    \item $(T,B)=(100000,5000)$ η σχετική διαφορά είναι $0.072664\%$ στο $L$, $0\%$ στο $P_{rej}$ και $0.032748\%$ στο $W$,
    \item $(T,B)=(1000000,20000)$ η σχετική διαφορά γίνεται $0.048009\%$, $0\%$ και $0.008514\%$ αντίστοιχα,
    \item $(T,B)=(10000000,120000)$ μειώνεται ακόμη περισσότερο σε $0.006206\%$, $0\%$ και $0.003304\%$.
\end{itemize}

Η εικόνα αυτή είναι απολύτως λογική. Όταν $\lambda=1$ και $\mu=5$, το σύστημα είναι σαφώς υποφορτισμένο, άρα ακόμη και αν ξεκινήσει άδειο, η δυναμική του δεν απέχει πολύ από τη μόνιμη κατάσταση. Με άλλα λόγια, η αρχική κατάσταση δεν στρεβλώνει ιδιαίτερα τις εκτιμήσεις, επειδή το σύστημα έτσι κι αλλιώς περνά το μεγαλύτερο μέρος του χρόνου σε χαμηλές καταστάσεις.

\subsubsection*{Περίπτωση $\lambda=5$}

Για $\lambda=5$ η επίδραση του transient είναι αισθητά εντονότερη, ιδιαίτερα στο μικρότερο run. Πιο συγκεκριμένα:
\begin{itemize}
    \item για $(T,B)=(100000,5000)$ οι σχετικές διαφορές είναι $0.792635\%$ στο $L$, $1.436403\%$ στο $P_{rej}$ και $0.902400\%$ στο $W$,
    \item για $(T,B)=(1000000,20000)$ μειώνονται σε $0.107573\%$, $0.299542\%$ και $0.128163\%$,
    \item για $(T,B)=(10000000,120000)$ γίνονται πλέον πολύ μικρές: $0.003459\%$, $0.019062\%$ και $0.008894\%$.
\end{itemize}

Η συγκεκριμένη συμπεριφορά είναι εξαιρετικά ενδιαφέρουσα. Στην περίπτωση $\lambda=\mu=5$, δηλαδή $\rho=1$, το σύστημα βρίσκεται ακριβώς στην οριακή περιοχή φόρτου. Αυτό σημαίνει ότι ούτε αδειάζει πολύ εύκολα ούτε οδηγείται αμέσως στην πλήρη κατάσταση. Ως αποτέλεσμα, η αρχική συνθήκη $n(0)=0$ επηρεάζει περισσότερο την εξέλιξη του συστήματος και χρειάζεται περισσότερος χρόνος μέχρι να απομακρυνθεί η επίδρασή της από τις εκτιμήσεις. Για τον λόγο αυτό, στην οριακή αυτή περίπτωση η αποκοπή του transient έχει τη μεγαλύτερη σημασία.

Αυτό επιβεβαιώνεται πλέον όχι μόνο από τις σχετικές διαφορές στα μεγέθη $L$, $P_{rej}$ και $W$, αλλά και από τη σύγκριση των ίδιων των εργοδικών πιθανοτήτων $P$ και $P_B$ στο ίδιο διάγραμμα, όπου η περίπτωση $\lambda=5$ εμφανίζει την πιο ορατή απόκλιση για μικρό $T$.

\subsubsection*{Περίπτωση $\lambda=10$}

Για $\lambda=10$ οι διαφορές είναι ξανά μικρές, αν και όχι μηδενικές:
\begin{itemize}
    \item για $(T,B)=(100000,5000)$ έχουμε $0.025674\%$ στο $L$, $0.012666\%$ στο $P_{rej}$ και $0.065978\%$ στο $W$,
    \item για $(T,B)=(1000000,20000)$ οι αντίστοιχες διαφορές είναι $0.005469\%$, $0.031798\%$ και $0.032481\%$,
    \item για $(T,B)=(10000000,120000)$ γίνονται $0.003702\%$, $0.011419\%$ και $0.012741\%$.
\end{itemize}

Και εδώ η εικόνα είναι αναμενόμενη. Το σύστημα ξεκινά πράγματι άδειο, όμως επειδή ο ρυθμός άφιξης είναι διπλάσιος από τον ρυθμό εξυπηρέτησης, τείνει πολύ γρήγορα να οδηγηθεί σε υψηλές καταστάσεις. Επομένως, παρότι υπάρχει transient, η μόνιμη συμπεριφορά επιβάλλεται σχετικά γρήγορα και έτσι η διαφορά μεταξύ με transient και χωρίς transient παραμένει μικρή.

\subsubsection*{Συμπέρασμα για το transient}

Συνολικά, η αφαίρεση της αρχικής μεταβατικής περιόδου επηρεάζει μεν τις εκτιμήσεις, αλλά η επίδρασή της μειώνεται συστηματικά όσο μεγαλώνει το $T$. Η επίδραση είναι πολύ μικρή για $\lambda=1$, αισθητή και πιο σημαντική για $\lambda=5$, και ξανά περιορισμένη για $\lambda=10$. Αυτό συμφωνεί με την εικόνα που δίνει και η θεωρία: το πόσο επικίνδυνη είναι η αρχική κατάσταση εξαρτάται από το πόσο κοντά βρίσκεται η δυναμική του συστήματος στο τεχνητό αρχικό σημείο $n(0)=0$.

\subsection{Μεταβολή της συμπεριφοράς του συστήματος καθώς αυξάνεται το λ}

Η δεύτερη βασική ερώτηση της άσκησης αφορά το πώς μεταβάλλεται η συμπεριφορά του συστήματος όταν αυξάνεται ο ρυθμός άφιξης. Η απάντηση προκύπτει τόσο από τα αριθμητικά αποτελέσματα όσο και από τα διαγράμματα εργοδικών πιθανοτήτων και σύγκλισης του $L$.

\subsubsection*{Για $\lambda=1$: σαφώς υποφορτισμένο σύστημα}

Στην περίπτωση αυτή έχουμε:
\[
\rho=\frac{\lambda}{\mu}=\frac{1}{5}=0.2.
\]
Το σύστημα είναι έντονα υποφορτισμένο και αυτό αποτυπώνεται καθαρά στα αποτελέσματα. Ο μέσος αριθμός πελατών παραμένει περίπου
\[
L \approx 0.25,
\]
ενώ ο μέσος χρόνος παραμονής είναι επίσης περίπου
\[
W \approx 0.25 \ \text{min}.
\]
Η πιθανότητα απόρριψης είναι στην πράξη μηδενική σε όλες τις εκτελέσεις.

Οι εργοδικές πιθανότητες δείχνουν ότι το σύστημα βρίσκεται συντριπτικά συχνότερα στην κατάσταση $0$, με
\[
P_0 \approx 0.8,
\]
ενώ οι πιθανότητες φθίνουν πολύ γρήγορα για τις επόμενες καταστάσεις:
\[
P_1 \approx 0.16,\qquad P_2 \approx 0.032,
\]
και στη συνέχεια γίνονται εξαιρετικά μικρές. Τόσο οι μπάρες του $P$ όσο και του $P_B$ ακολουθούν αυτή την εικόνα σχεδόν ταυτιζόμενες. Άρα, όταν ο ρυθμός άφιξης είναι αρκετά μικρότερος από τον ρυθμό εξυπηρέτησης, το σύστημα λειτουργεί ομαλά, με ελάχιστη συμφόρηση και πρακτικά μηδενικές απώλειες.

\subsubsection*{Για $\lambda=5$: οριακή κατάσταση φόρτου}

Εδώ έχουμε:
\[
\rho=\frac{5}{5}=1.
\]
Αυτή είναι η πιο ενδιαφέρουσα περίπτωση, διότι βρίσκεται ακριβώς στο όριο μεταξύ υποφορτισμένης και υπερφορτισμένης συμπεριφοράς.

Τα αποτελέσματα δείχνουν ότι:
\[
L \approx 5,\qquad P_{rej}\approx 0.091,\qquad W\approx 1.10 \ \text{min}.
\]
Η τιμή της πιθανότητας απόρριψης είναι κοντά στη θεωρητική τιμή:
\[
P_{rej}=P_{10}=\frac{1}{11}\approx 0.090909.
\]

Τα διαγράμματα των εργοδικών πιθανοτήτων είναι επίσης πολύ χαρακτηριστικά: όλες οι πιθανότητες είναι σχεδόν ίσες μεταξύ τους, δηλαδή περίπου
\[
P_0\approx P_1\approx \cdots \approx P_{10}\approx \frac{1}{11}.
\]
Αυτό ακριβώς είναι το αναμενόμενο θεωρητικό αποτέλεσμα για την περίπτωση $\rho=1$ σε σύστημα $\mathrm{M/M/1/K}$.

Η ταυτόχρονη σχεδίαση των $P$ και $P_B$ δείχνει ότι στο μικρότερο run υπάρχει μία μικρή αλλά ορατή απόκλιση μεταξύ των δύο κατανομών, η οποία όμως μειώνεται πολύ γρήγορα καθώς αυξάνεται το $T$. Άρα, στην οριακή αυτή περιοχή, το σύστημα επισκέπτεται με παρόμοια συχνότητα όλες τις καταστάσεις του, από την άδεια μέχρι και την πλήρη. Αυτό σημαίνει ότι η συμφόρηση είναι πλέον ουσιώδης, ενώ η πιθανότητα blocking δεν είναι αμελητέα, αλλά παραμένει ακόμη σε σχετικά μέτρια επίπεδα.

\subsubsection*{Για $\lambda=10$: έντονα φορτισμένο σύστημα}

Στην περίπτωση αυτή:
\[
\rho=\frac{10}{5}=2.
\]
Το σύστημα δέχεται αφίξεις με ρυθμό διπλάσιο από αυτόν που μπορεί να εξυπηρετήσει. Αν επρόκειτο για άπειρη ουρά, η κατάσταση αυτή θα οδηγούσε σε μη φραγμένη αύξηση του πλήθους των πελατών. Εδώ όμως το σύστημα είναι πεπερασμένο και διαθέτει χωρητικότητα μόνο $10$ θέσεων. Επομένως, αντί να αυξάνει απεριόριστα, κολλάει συχνά κοντά στη μέγιστη χωρητικότητα και απορρίπτει σημαντικό ποσοστό των νέων αφίξεων.

Πράγματι, στα αποτελέσματα βλέπουμε:
\[
L \approx 9.0,\qquad P_{rej}\approx 0.50,\qquad W\approx 1.80 \ \text{min}.
\]
Ιδιαίτερα εντυπωσιακό είναι ότι η πιθανότητα η κατάσταση να είναι ακριβώς $10$ είναι περίπου:
\[
P_{10}\approx 0.50,
\]
δηλαδή περίπου το μισό του χρόνου το σύστημα είναι πλήρες. Εξίσου χαρακτηριστικό είναι ότι οι πιθανότητες των καταστάσεων αυξάνουν καθώς πηγαίνουμε από χαμηλά προς υψηλά state, με τις μεγαλύτερες μάζες να συγκεντρώνονται στις καταστάσεις $8$, $9$ και $10$.

Οι δύο κατανομές $P$ και $P_B$ είναι και εδώ πολύ κοντά μεταξύ τους, πράγμα που δείχνει ότι η επίδραση του transient στην έντονα φορτισμένη περίπτωση είναι μικρότερη από ό,τι στην οριακή περίπτωση $\lambda=5$.

Άρα, καθώς το $\lambda$ αυξάνεται, το σύστημα μεταβαίνει από μία περιοχή ομαλής και άνετης λειτουργίας, σε μία οριακή κατάσταση εξισορρόπησης, και τελικά σε μία κατάσταση έντονης συμφόρησης με ισχυρό blocking. Η μετάβαση αυτή αποτυπώνεται ταυτόχρονα:
\begin{itemize}
    \item στο σχήμα των εργοδικών πιθανοτήτων,
    \item στην αύξηση του $L$,
    \item στην αύξηση του $P_{rej}$,
    \item και στην αύξηση του $W$.
\end{itemize}

\subsection{Σύγκριση των αποτελεσμάτων με τη θεωρία}

Ένας πολύ ισχυρός τρόπος επαλήθευσης της ορθότητας της προσομοίωσης είναι η σύγκριση των αριθμητικών αποτελεσμάτων με τις θεωρητικές τιμές της ουράς $\mathrm{M/M/1/10}$.

\subsubsection*{Θεωρητικές τιμές για $\lambda=1$}

Για $\lambda=1$ και $\mu=5$ έχουμε $\rho=0.2$. Θεωρητικά:
\[
P_0 \approx 0.800000,\qquad P_1 \approx 0.160000,\qquad P_2 \approx 0.032000,
\]
\[
L \approx 0.250000,\qquad W \approx 0.250000 \ \text{min},
\]
ενώ η πιθανότητα απόρριψης είναι:
\[
P_{rej}=P_{10}\approx 8.19\times 10^{-8},
\]
δηλαδή ουσιαστικά μηδενική για τους πρακτικούς σκοπούς της άσκησης.

Η προσομοίωση έδωσε, στην καλύτερη εκτέλεση,
\[
L=0.2504273515,\qquad L_B=0.2504118097,
\]
\[
W=0.2502386518,\qquad W_B=0.2502303841,
\]
και πρακτικά μηδενικά $P_{rej}$ και $P_{rej,B}$. Η συμφωνία είναι εξαιρετική. Το ίδιο ισχύει και για τις εργοδικές πιθανότητες, όπου οι δύο εκτιμήσεις $P$ και $P_B$ σχεδόν συμπίπτουν με την αναμενόμενη γεωμετρικά φθίνουσα μορφή της κατανομής.

\subsubsection*{Θεωρητικές τιμές για $\lambda=5$}

Για $\lambda=5$ έχουμε $\rho=1$, άρα θεωρητικά:
\[
P_n=\frac{1}{11}\approx 0.090909,\qquad n=0,1,\ldots,10.
\]
Συνεπώς:
\[
L=\frac{0+1+\cdots+10}{11}=5,
\]
\[
P_{rej}=P_{10}=\frac{1}{11}\approx 0.090909,
\]
και
\[
\lambda_{\text{eff}}=5\left(1-\frac{1}{11}\right)=\frac{50}{11}\approx 4.54545,
\]
οπότε:
\[
W=\frac{L}{\lambda_{\text{eff}}}=\frac{5}{50/11}=1.1 \ \text{min}.
\]

Η προσομοίωση έδωσε:
\[
L=5.0076954577,\qquad L_B=5.0075222213,
\]
\[
P_{rej}=0.0911728040,\qquad P_{rej,B}=0.0911554245,
\]
\[
W=1.1017529664,\qquad W_B=1.1016549734.
\]
Και εδώ η συμφωνία με τη θεωρία είναι πολύ καλή.

\subsubsection*{Θεωρητικές τιμές για $\lambda=10$}

Για $\lambda=10$ και $\mu=5$ έχουμε $\rho=2$. Θεωρητικά, οι πιθανότητες είναι ανάλογες των δυνάμεων του $2$, άρα οι μεγαλύτερες καταστάσεις αποκτούν τη συντριπτική πλειονότητα της μάζας. Ειδικά:
\[
P_{10}\approx 0.500244,\qquad P_9\approx 0.250122,\qquad P_8\approx 0.125061.
\]
Ακόμη:
\[
L\approx 9.005374,\qquad P_{rej}\approx 0.500244,\qquad W\approx 1.801955 \ \text{min}.
\]

Η προσομοίωση στην καλύτερη εκτέλεση έδωσε:
\[
L=9.0065989686,\qquad L_B=9.0062655890,
\]
\[
P_{rej}=0.5004547366,\qquad P_{rej,B}=0.5003975920,
\]
\[
W=1.8020232187,\qquad W_B=1.8017936156,
\]
καθώς και:
\[
P_{10}\approx 0.500292,\qquad P_{10,B}\approx 0.500233.
\]
Η απόκλιση από τη θεωρία είναι ξανά πολύ μικρή, γεγονός που δείχνει ότι η προσομοίωση είναι όχι μόνο λογικά σωστή, αλλά και αριθμητικά αξιόπιστη.

\subsection{Σύγκλιση των εκτιμήσεων και επίδραση του $T$}

Ζητείται ρητά να συγκρίνουμε τα διαφορετικά ζεύγη $(T,B)$ ως προς τη σύγκλιση. Η επίδραση του $T$ φαίνεται τόσο από τα αριθμητικά αποτελέσματα όσο και από τα διαγράμματα εξέλιξης του $L$.

Συγκεκριμένα, όσο αυξάνεται ο αριθμός μεταβάσεων:
\[
10^5 \to 10^6 \to 10^7,
\]
παρατηρούμε ότι:

\begin{itemize}
    \item οι σχετικές διαφορές μεταξύ με transient και χωρίς transient μειώνονται συστηματικά,
    \item οι εκτιμήσεις των μεγεθών σταθεροποιούνται,
    \item οι καμπύλες του $L$ ως προς τον αριθμό μεταβάσεων εμφανίζουν μικρότερες τυχαίες διακυμάνσεις σε σχέση με το συνολικό μήκος της εκτέλεσης,
    \item οι δύο κατανομές εργοδικών πιθανοτήτων $P$ και $P_B$ πλησιάζουν όλο και περισσότερο η μία την άλλη,
    \item τα αποτελέσματα πλησιάζουν όλο και περισσότερο τις θεωρητικές τιμές.
\end{itemize}

Με άλλα λόγια, το $T$ λειτουργεί ακριβώς όπως θα αναμέναμε σε μία μη τερματική προσομοίωση: όσο μεγαλύτερη είναι η διάρκεια της εκτέλεσης, τόσο μικρότερη γίνεται η στατιστική αβεβαιότητα των εκτιμήσεων.

Αυτό φαίνεται ιδιαίτερα καθαρά στην περίπτωση $\lambda=5$, όπου το μικρότερο run εμφανίζει την πιο αισθητή διαφορά μεταξύ των δύο τρόπων υπολογισμού, αλλά και τη μεγαλύτερη ορατή διαφορά μεταξύ των $P$ και $P_B$ στο grouped bar plot, ενώ στο μεγαλύτερο run οι αποκλίσεις έχουν σχεδόν μηδενιστεί. Η οριακή αυτή περίπτωση είναι, επομένως, και η πιο απαιτητική ως προς τη διάρκεια της προσομοίωσης.

\subsection{Επάρκεια της επιλογής του $B$ ως transient period}

Ένα ακόμη ζητούμενο της άσκησης είναι να αξιολογηθεί αν οι τιμές:
\[
B=5000,\qquad B=20000,\qquad B=120000
\]
είναι επαρκείς ως αποκοπή της αρχικής μεταβατικής περιόδου για τα αντίστοιχα $T$.

Με βάση τα αποτελέσματα, μπορούμε να πούμε τα εξής:

\begin{itemize}
    \item Για $\lambda=1$, οι τιμές του $B$ είναι σαφώς επαρκείς. Η επίδραση της αρχικής κατάστασης είναι εξαρχής πολύ μικρή και οι διαφορές μεταξύ των δύο εκτιμήσεων είναι αμελητέες.
    \item Για $\lambda=10$, οι τιμές του $B$ είναι επίσης επαρκείς. Παρότι το σύστημα ξεκινά άδειο, μετακινείται πολύ γρήγορα προς τις υψηλές καταστάσεις και άρα \say{ξεχνά} σχετικά σύντομα την αρχική συνθήκη.
    \item Για $\lambda=5$, η επιλογή του $B$ είναι η πιο κρίσιμη. Στο μικρότερο run, η επίδραση του transient παραμένει εμφανής, άρα θα μπορούσε κανείς, αν επιζητούσε ακόμη πιο καθαρό διαχωρισμό, να προτείνει ένα ελαφρώς μεγαλύτερο $B$. Ωστόσο, στα μεγαλύτερα runs η τρέχουσα επιλογή του $B$ αποδεικνύεται επαρκής, αφού οι σχετικές διαφορές μειώνονται σε πολύ χαμηλά επίπεδα και οι κατανομές $P$ και $P_B$ πρακτικά συμπίπτουν.
\end{itemize}

Συνεπώς, δεν προκύπτει κάποια σοβαρή ένδειξη ότι οι δοθείσες τιμές του $B$ είναι λανθασμένες. Αντίθετα, φαίνεται ότι αποτελούν μία αρκετά λογική και ισορροπημένη επιλογή για τα ζητούμενα μεγέθη της άσκησης. Μόνο στην οριακή περίπτωση $\lambda=5$ και στο μικρότερο run θα μπορούσε να υποστηριχθεί ότι ένα λίγο μεγαλύτερο transient cutoff ίσως έδινε ελαφρώς ακόμη πιο καθαρές εκτιμήσεις.

\subsection{Σύγκριση των χρόνων εκτέλεσης}

Παραθέτουμε συνοπτικά τους χρόνους εκτέλεσης των $9$ runs:

\begin{table}[h!]
\centering
\caption{Χρόνοι εκτέλεσης της προσομοίωσης.}
\begin{tabular}{|c|c|c|c|}
\hline
$\lambda$ & $(10^5,5000)$ & $(10^6,20000)$ & $(10^7,120000)$ \\ \hline
$1$  & $6.6001$ s  & $65.9007$ s & $635.6335$ s \\ \hline
$5$  & $6.9150$ s  & $67.4362$ s & $639.4556$ s \\ \hline
$10$ & $5.8064$ s  & $59.6151$ s & $626.7589$ s \\ \hline
\end{tabular}
\end{table}

Ο βασικός παράγοντας που καθορίζει τον χρόνο εκτέλεσης είναι προφανώς ο συνολικός αριθμός μεταβάσεων $T$. Όταν περνάμε από:
\[
10^5 \to 10^6 \to 10^7,
\]
ο χρόνος εκτέλεσης αυξάνεται πολύ σημαντικά, όπως είναι απολύτως αναμενόμενο.

Αντίθετα, η εξάρτηση του χρόνου εκτέλεσης από το $\lambda$ δεν είναι απολύτως μονοτονική στα δικά μας αποτελέσματα. Αυτό είναι λογικό, διότι στη συγκεκριμένη υλοποίηση παίζουν ρόλο και επιμέρους πρακτικές λεπτομέρειες της εκτέλεσης στο Octave, όπως η ακριβής ακολουθία γεγονότων, οι προσπελάσεις στους πίνακες και γενικότερα η υλοποιητική επιβάρυνση της κάθε διαδρομής του κώδικα. Άρα, το ουσιώδες συμπέρασμα δεν είναι η μικρή μεταβολή του χρόνου με το $\lambda$, αλλά η πολύ ισχυρή αύξηση του χρόνου εκτέλεσης με το $T$.

\subsection{Σύνθεση των παρατηρήσεων ανά τιμή του λ}

Αν συγκεντρώσουμε όλες τις παρατηρήσεις σε μία ενιαία εικόνα, προκύπτουν τα εξής:

\paragraph{Για $\lambda=1$.}
Το σύστημα λειτουργεί άνετα. Η πιθανότητα να είναι άδειο είναι πολύ μεγάλη, το blocking είναι πρακτικά ανύπαρκτο, ο μέσος αριθμός πελατών είναι πολύ μικρός και ο μέσος χρόνος παραμονής παραμένει κοντά στο $0.25$ min. Η αφαίρεση του transient επηρεάζει ελάχιστα τα αποτελέσματα και οι δύο εκδοχές των εργοδικών πιθανοτήτων είναι σχεδόν ίδιες.

\paragraph{Για $\lambda=5$.}
Το σύστημα βρίσκεται σε οριακή περιοχή. Όλες οι καταστάσεις εμφανίζονται σχεδόν με την ίδια πιθανότητα, ο μέσος αριθμός πελατών είναι περίπου $5$, η πιθανότητα blocking είναι περίπου $9.1\%$ και ο μέσος χρόνος παραμονής περίπου $1.1$ min. Εδώ η αφαίρεση του transient έχει τη μεγαλύτερη σημασία και το φαινόμενο της σύγκλισης με αύξηση του $T$ φαίνεται καθαρότερα τόσο στα μεγέθη όσο και στα ίδια τα bar plots.

\paragraph{Για $\lambda=10$.}
Το σύστημα βρίσκεται σε έντονα φορτισμένη λειτουργία. Παραμένει για μεγάλο μέρος του χρόνου στις καταστάσεις $9$ και $10$, η πιθανότητα blocking γίνεται περίπου $50\%$, και ο μέσος αριθμός πελατών πλησιάζει το μέγιστο θεωρητικά εφικτό επίπεδο για το συγκεκριμένο capacity. Το transient υπάρχει αλλά αποσβένεται σχετικά γρήγορα, ενώ οι δύο εκδοχές των εργοδικών πιθανοτήτων διαφέρουν ελάχιστα.

\subsection{Άμεσες απαντήσεις στα ερωτήματα της άσκησης}

Με βάση τα παραπάνω, οι ζητούμενες απαντήσεις της εκφώνησης συνοψίζονται ως εξής:

\begin{enumerate}[label=\arabic*.]
    \item \textbf{Πώς επηρεάζει η αφαίρεση της αρχικής μεταβατικής περιόδου τις εκτιμήσεις των μεγεθών;}\\
    Η αφαίρεση του transient βελτιώνει τις εκτιμήσεις, απομακρύνοντας τη μεροληψία που προκαλεί η τεχνητή αρχική κατάσταση $n(0)=0$. Η επίδρασή της είναι πολύ μικρή για $\lambda=1$, μεγαλύτερη για $\lambda=5$ και ξανά σχετικά μικρή για $\lambda=10$. Σε όλες τις περιπτώσεις, καθώς αυξάνει το $T$, η διαφορά ανάμεσα στις εκτιμήσεις με και χωρίς transient μειώνεται.

    \item \textbf{Πώς μεταβάλλεται η συμπεριφορά του συστήματος καθώς αυξάνεται η τιμή του $\lambda$;}\\
    Όσο αυξάνει το $\lambda$, το σύστημα μεταβαίνει από ομαλή λειτουργία με μικρό πλήθος πελατών και ουσιαστικά μηδενικό blocking, σε οριακή λειτουργία με περίπου ομοιόμορφη κατανομή καταστάσεων, και τελικά σε έντονα συμφραζόμενη λειτουργία όπου το σύστημα μένει συχνά σχεδόν πλήρες και απορρίπτει περίπου τις μισές αφίξεις.

    \item \textbf{Αν η επιλογή $B$ αρχικών μεταβάσεων σε κάθε ζεύγος ως transient period φαίνεται επαρκής ή αν θα προτείνατε διαφορετική τιμή;}\\
    Η επιλογή των δοθέντων τιμών του $B$ φαίνεται γενικά επαρκής. Για $\lambda=1$ και $\lambda=10$ είναι σαφώς ικανοποιητική. Για $\lambda=5$ και ειδικά στο μικρότερο run θα μπορούσε ίσως να εξεταστεί ένα ελαφρώς μεγαλύτερο $B$, αλλά τα αποτελέσματα δείχνουν ότι ακόμη και η τρέχουσα επιλογή παραμένει εύλογη και λειτουργική.

    \item \textbf{Σύγκριση των διαφορετικών τιμών των $T$ και $B$ ως προς σύγκλιση, χρόνο εκτέλεσης και σχετική διαφορά με και χωρίς transient period.}\\
    Μεγαλύτερα $T$ οδηγούν σε καλύτερη σύγκλιση, μικρότερες σχετικές διαφορές και αποτελέσματα που πλησιάζουν περισσότερο τη θεωρία, με τίμημα όμως πολύ μεγαλύτερο χρόνο εκτέλεσης. Επιπλέον, η σύγκριση των $P$ και $P_B$ στο ίδιο γράφημα δείχνει οπτικά την ίδια τάση: όσο μεγαλώνει το $T$, τόσο οι δύο κατανομές τείνουν να συμπέσουν. Το κύριο συμπέρασμα είναι ότι η αύξηση του $T$ βελτιώνει ουσιαστικά την αξιοπιστία των εκτιμήσεων, ιδιαίτερα στην οριακή περίπτωση $\lambda=5$.
\end{enumerate}

\subsection{Τελικό συμπέρασμα}

Η προσομοίωση του συστήματος $\mathrm{M/M/1/10}$ παρήγαγε αποτελέσματα απολύτως συμβατά τόσο με τη θεωρία όσο και με τη φυσική ερμηνεία του μοντέλου. Η αύξηση του ρυθμού αφίξεων οδηγεί σε αύξηση της συμφόρησης, του μέσου αριθμού πελατών και της πιθανότητας απόρριψης, ενώ η αφαίρεση της αρχικής μεταβατικής περιόδου συμβάλλει στη βελτίωση των εκτιμήσεων, ιδίως όταν το σύστημα βρίσκεται σε οριακή κατάσταση φόρτου.

Επιπλέον, τα αποτελέσματα για τα μεγαλύτερα $T$ εμφανίζουν πολύ καλή σταθερότητα και πολύ μικρή διαφορά μεταξύ με transient και χωρίς transient, τόσο στα αριθμητικά μεγέθη όσο και στις ίδιες τις εργοδικές πιθανότητες. Η μορφή των διαγραμμάτων καθιστά αυτή τη σύγκριση άμεση και οπτικά σαφή. Συνεπώς, θεωρούμε ότι τόσο η υλοποίηση όσο και τα παραγόμενα αριθμητικά αποτελέσματα είναι επιτυχή και αξιόπιστα.

\newpage
\section{Τροποποίηση του μοντέλου για ρυθμό εξυπηρέτησης που εξαρτάται από την κατάσταση}

\subsection{Νέα διατύπωση του μοντέλου}

Στο προηγούμενο μέρος της άσκησης θεωρήσαμε ότι ο ρυθμός εξυπηρέτησης ήταν σταθερός και ίσος με
\[
\mu = 5 \ \text{πελάτες/min},
\]
ανεξάρτητα από την κατάσταση του συστήματος. Στο παρόν ερώτημα, η εκφώνηση ζητά να περιγράψουμε πώς θα άλλαζε η προσομοίωση αν ο ρυθμός εξυπηρέτησης δεν ήταν σταθερός, αλλά εξαρτιόταν από την τρέχουσα κατάσταση του συστήματος. Συγκεκριμένα, στην κατάσταση $i$ ο ρυθμός εξυπηρέτησης δίνεται από:
\[
\mu_i = \mu(i+1), \qquad i=1,2,\ldots,10,
\]
όπου
\[
\mu = 1 \ \text{πελάτης/sec}.
\]

Άρα, όταν το σύστημα έχει:
\[
i=1 \Rightarrow \mu_1 = 2 \ \text{πελάτες/sec},
\]
\[
i=2 \Rightarrow \mu_2 = 3 \ \text{πελάτες/sec},
\]
\[
\cdots
\]
\[
i=10 \Rightarrow \mu_{10} = 11 \ \text{πελάτες/sec}.
\]

Με άλλα λόγια, όσο αυξάνεται ο αριθμός πελατών μέσα στο σύστημα, τόσο αυξάνεται και ο ρυθμός με τον οποίο ολοκληρώνονται οι εξυπηρετήσεις. Επομένως, το μοντέλο παραμένει μεν σύστημα $\mathrm{M/M/1/10}$ ως προς τη χωρητικότητα και τη Markovian φύση των αφίξεων και εξυπηρετήσεων, όμως παύει να είναι το κλασικό $\mathrm{M/M/1/K}$ με σταθερό death rate. Πλέον πρόκειται για birth--death process με:

\[
\lambda_i =
\begin{cases}
\lambda, & i=0,1,\ldots,9,\\
0, & i=10,
\end{cases}
\]
και
\[
\mu_i =
\begin{cases}
0, & i=0,\\
\mu(i+1), & i=1,2,\ldots,10.
\end{cases}
\]

Άρα το state space παραμένει:
\[
\{0,1,2,\ldots,10\},
\]
όμως οι μεταβάσεις προς τα κάτω γίνονται πλέον ταχύτερες όσο το σύστημα φορτώνεται περισσότερο.

\subsection{Αλλαγή στις μονάδες μέτρησης}

Στο σημείο αυτό απαιτείται ιδιαίτερη προσοχή στις μονάδες. Στα προηγούμενα ερωτήματα χρησιμοποιήσαμε:
\[
\lambda \in \{1,5,10\} \ \text{πελάτες/min}
\]
και
\[
\mu = 5 \ \text{πελάτες/min}.
\]
Στο παρόν ερώτημα, η νέα παράμετρος δίνεται ως:
\[
\mu = 1 \ \text{πελάτης/sec}.
\]

Επομένως, πριν από οποιαδήποτε αλλαγή στον κώδικα, πρέπει υποχρεωτικά να ενοποιηθούν οι μονάδες του χρόνου. Υπάρχουν δύο ισοδύναμοι τρόποι:

\begin{enumerate}[label=\arabic*.]
    \item είτε κρατάμε όλη την προσομοίωση σε λεπτά και μετατρέπουμε:
    \[
    1 \ \text{πελάτης/sec} = 60 \ \text{πελάτες/min},
    \]
    οπότε
    \[
    \mu_i = 60(i+1) \ \text{πελάτες/min},
    \]
    \item είτε κρατάμε όλη την προσομοίωση σε δευτερόλεπτα και μετατρέπουμε τους ρυθμούς αφίξεων:
    \[
    \lambda \in \left\{\frac{1}{60},\frac{5}{60},\frac{10}{60}\right\} \ \text{πελάτες/sec}.
    \]
\end{enumerate}

Στη δική μας υλοποίηση, επειδή όλος ο προηγούμενος κώδικας ήταν γραμμένος με χρόνους σε λεπτά, η φυσικότερη επιλογή θα ήταν να κρατήσουμε ως μονάδα τα λεπτά και να μετατρέψουμε τον νέο βασικό ρυθμό εξυπηρέτησης σε:
\[
\mu_{\text{base}} = 60 \ \text{πελάτες/min}.
\]
Τότε, στην κατάσταση $i$,
\[
\mu_i = 60(i+1) \ \text{πελάτες/min}.
\]

Αν δεν γίνει αυτή η μετατροπή και χρησιμοποιηθούν ταυτόχρονα ρυθμοί αφίξεων σε πελάτες/min και ρυθμοί εξυπηρέτησης σε πελάτες/sec, τότε η προσομοίωση θα είναι μαθηματικά ασυνεπής.

\subsection{Πώς αλλάζει η παραγωγή των χρόνων εξυπηρέτησης}

Στο αρχικό μοντέλο με σταθερό ρυθμό εξυπηρέτησης, κάθε χρόνος εξυπηρέτησης παραγόταν από εκθετική κατανομή με σταθερό ρυθμό:
\[
S \sim \mathrm{Exp}(\mu).
\]
Αυτό σήμαινε ότι, κάθε φορά που το σύστημα ήταν μη κενό, ο μηχανισμός αναχωρήσεων είχε τον ίδιο ακριβώς ρυθμό, ανεξάρτητα από το αν υπήρχαν $1$, $5$ ή $10$ πελάτες στο σύστημα.

Στο νέο μοντέλο αυτό δεν ισχύει πλέον. Όταν το σύστημα βρίσκεται στην κατάσταση $i>0$, ο χρόνος μέχρι την επόμενη αναχώρηση πρέπει να παράγεται από:
\[
S_i \sim \mathrm{Exp}(\mu_i)=\mathrm{Exp}(\mu(i+1)).
\]
Άρα ο χρόνος εξυπηρέτησης δεν εξαρτάται απλώς από μια σταθερή παράμετρο, αλλά από το τρέχον state.

Συνεπώς:

\begin{itemize}
    \item σε μικρές καταστάσεις ο ρυθμός αναχώρησης είναι σχετικά μικρός,
    \item σε μεγάλες καταστάσεις ο ρυθμός αναχώρησης είναι μεγαλύτερος,
    \item και επομένως η πιθανότητα να συμβεί σύντομα αναχώρηση αυξάνεται όσο μεγαλώνει το πλήθος πελατών στο σύστημα.
\end{itemize}

Αυτό είναι ακριβώς και το βασικό ποιοτικό συμπέρασμα της τροποποίησης: το σύστημα αποκτά έναν μηχανισμό αυτο-ενίσχυσης της εξυπηρέτησης όταν αυξάνεται η συμφόρηση.

\subsection{Πώς αλλάζει ο μηχανισμός επιλογής του επόμενου γεγονότος}

Η δική μας υλοποίηση δεν στηρίζεται σε ξεχωριστή εκ των προτέρων παραγωγή ενός πλήρους service time για κάθε πελάτη, αλλά σε \foreignlanguage{english}{event-driven} λογική: σε κάθε μετάβαση παράγεται ο χρόνος μέχρι το επόμενο event και στη συνέχεια αποφασίζεται αν αυτό είναι άφιξη ή αναχώρηση.

Στην αρχική εκδοχή του κώδικα, όταν το σύστημα ήταν μη κενό, είχαμε:
\[
\text{total\_rate} = \lambda + \mu
\]
και:
\[
P(\text{arrival} \mid state>0)=\frac{\lambda}{\lambda+\mu},
\qquad
P(\text{departure} \mid state>0)=\frac{\mu}{\lambda+\mu}.
\]

Με τον νέο state-dependent ρυθμό, τα παραπάνω γίνονται:
\[
\text{total\_rate}(i)=\lambda+\mu_i=\lambda+\mu(i+1), \qquad i=1,\ldots,10,
\]
και άρα:
\[
P(\text{arrival} \mid state=i)=\frac{\lambda}{\lambda+\mu(i+1)},
\]
\[
P(\text{departure} \mid state=i)=\frac{\mu(i+1)}{\lambda+\mu(i+1)}.
\]

Επιπλέον, ο χρόνος μέχρι το επόμενο γεγονός πρέπει να δειγματοληπτείται από:
\[
\Delta t \sim \mathrm{Exp}\bigl(\lambda+\mu(i+1)\bigr), \qquad i>0.
\]

Αντίθετα, όταν το σύστημα είναι άδειο, δηλαδή όταν
\[
state=0,
\]
δεν υπάρχει μηχανισμός departure, οπότε παραμένει:
\[
\text{total\_rate} = \lambda,
\]
και το επόμενο γεγονός είναι υποχρεωτικά άφιξη.

\subsection{Ποια σημεία του κώδικα θα έπρεπε να τροποποιηθούν}

Σε επίπεδο υλοποίησης, οι απαραίτητες αλλαγές εντοπίζονται σε συγκεκριμένα σημεία.

\subsubsection*{1. Ορισμός της παραμέτρου εξυπηρέτησης}

Στην αρχική εκδοχή του προγράμματος είχαμε μία σταθερή τιμή:
\[
\mu = 5.
\]
Στη νέα εκδοχή δεν αρκεί μία μόνο σταθερά για όλες τις καταστάσεις. Θα πρέπει να οριστεί μία βασική τιμή, π.χ.
\[
\mu_{\text{base}} = 60 \ \text{πελάτες/min}
\]
αν επιλέξουμε να δουλέψουμε σε λεπτά, και στη συνέχεια σε κάθε βήμα να υπολογίζεται:
\[
\mu_{\text{state}} = \mu_{\text{base}} \cdot (state+1), \qquad \text{όταν } state>0.
\]

\subsubsection*{2. Τμήμα επιλογής event μέσα στο βασικό loop}

Στον αρχικό κώδικα, το κρίσιμο κομμάτι ήταν της μορφής:
\[
\texttt{total\_rate = lambda + mu}
\]
και
\[
\texttt{is\_arrival = (rand() < lambda / (lambda + mu))}.
\]

Αυτό θα έπρεπε να αντικατασταθεί από state-dependent υπολογισμό. Δηλαδή, όταν $state>0$:
\[
\mu_{\text{state}} = \mu_{\text{base}}(state+1),
\]
\[
\texttt{total\_rate = lambda + mu\_state},
\]
\[
\texttt{is\_arrival = (rand() < lambda / total\_rate)}.
\]

Άρα, η βασική λογική του loop δεν αλλάζει ως προς τη δομή, αλλά αλλάζει το γεγονός ότι ο ρυθμός departure δεν είναι πλέον σταθερός.

\subsubsection*{3. Παραγωγή του χρόνου μέχρι το επόμενο event}

Στην αρχική μορφή, για μη κενό σύστημα είχαμε:
\[
\Delta t = -\frac{\ln U}{\lambda+\mu}.
\]
Τώρα θα έχουμε:
\[
\Delta t = -\frac{\ln U}{\lambda+\mu_{\text{state}}}.
\]

Άρα το σημείο στο οποίο ο κώδικας δειγματοληπτεί το \texttt{dt} πρέπει επίσης να συνδεθεί με το νέο state-dependent \texttt{mu\_state}.

\subsubsection*{4. Επανυπολογισμός του departure rate μετά από κάθε αλλαγή state}

Η εισαγωγή της άσκησης αναφέρει ρητά ότι, όταν ο ρυθμός αναχώρησης εξαρτάται από την κατάσταση, μετά από κάθε άφιξη ή αναχώρηση που αλλάζει το $n$, πρέπει να επανυπολογίζεται το departure clock με βάση το νέο state. Αυτό σημαίνει ότι μετά από κάθε μετάβαση:
\[
state \to state_{\text{new}},
\]
στην επόμενη επανάληψη του loop πρέπει να χρησιμοποιείται:
\[
\mu_{\text{state\_new}} = \mu_{\text{base}}(state_{\text{new}}+1),
\]
εφόσον βέβαια
\[
state_{\text{new}} > 0.
\]

Στη δική μας υλοποίηση αυτό γίνεται φυσικά, επειδή σε κάθε νέο βήμα ο μηχανισμός event selection ξαναχτίζεται από το τρέχον state. Δηλαδή, δεν κουβαλάμε σταθερό departure rate από το παρελθόν, αλλά το υπολογίζουμε εκ νέου σε κάθε iteration.

\subsubsection*{5. Τα υπόλοιπα τμήματα του κώδικα που παραμένουν ίδια}

Είναι πολύ σημαντικό να τονιστεί ότι η υπόλοιπη δομή της προσομοίωσης δεν χρειάζεται ουσιαστική αλλαγή. Συγκεκριμένα, παραμένουν ακριβώς ίδια:

\begin{itemize}
    \item το state space $\{0,\ldots,10\}$,
    \item ο κανόνας blocking όταν $state=10$,
    \item ο μετρητής προσπαθειών άφιξης,
    \item ο μετρητής απορρίψεων,
    \item η FIFO διαχείριση των χρονικών στιγμών άφιξης για τον υπολογισμό του $W$,
    \item ο υπολογισμός των εργοδικών πιθανοτήτων μέσω των time-average counters,
    \item ο υπολογισμός των $L$, $P_{rej}$ και $W$,
    \item η ξεχωριστή συλλογή στατιστικών με και χωρίς transient period,
    \item η χρήση ίδιου seed σε όλα τα runs,
    \item το debugging trace των πρώτων μεταβάσεων.
\end{itemize}

Άρα η τροποποίηση είναι τοπική ως προς τη δυναμική του μηχανισμού departure, αλλά όχι ως προς τη συνολική αρχιτεκτονική της προσομοίωσης.

\subsection{Εναλλακτική ερμηνεία μέσω ανεξάρτητων clocks}

Αν αντί για τη δική μας συμπαγή \foreignlanguage{english}{event-driven} υλοποίηση χρησιμοποιούσαμε την κλασική λογική δύο ανεξάρτητων ρολογιών, τότε η τροποποίηση θα μπορούσε να περιγραφεί ως εξής:

\begin{itemize}
    \item το επόμενο interarrival time θα παρέμενε:
    \[
    A \sim \mathrm{Exp}(\lambda),
    \]
    \item ενώ το επόμενο service completion time θα γινόταν:
    \[
    D_i \sim \mathrm{Exp}(\mu_i)=\mathrm{Exp}(\mu(i+1)),
    \]
    όταν το σύστημα είναι στην κατάσταση $i>0$.
\end{itemize}

Τότε, σε κάθε χρονική στιγμή θα συγκρίναμε το \texttt{nextArrival} και το \texttt{nextDeparture} και θα εκτελούσαμε το γεγονός με τον μικρότερο χρόνο. Μετά από κάθε μεταβολή του state, το departure clock θα έπρεπε να επαναδημιουργείται με τον νέο ρυθμό. Αυτή είναι ακριβώς η λογική που περιγράφεται και στις διαφάνειες της άσκησης.

Η δική μας υλοποίηση είναι ισοδύναμη, αλλά πιο συμπαγής: δεν κρατά δύο ξεχωριστά clocks, αλλά χρησιμοποιεί το άθροισμα των ρυθμών και επιλέγει τύπο γεγονότος με την κατάλληλη υπό συνθήκη πιθανότητα.

\subsection{Αναμενόμενη επίδραση στα αποτελέσματα}

Αν εφαρμόζαμε πράγματι το νέο μοντέλο, θα αναμέναμε ουσιώδη αλλαγή στη συμπεριφορά του συστήματος σε σχέση με τα αποτελέσματα που είδαμε στα προηγούμενα ερωτήματα.

Εφόσον ο ρυθμός εξυπηρέτησης αυξάνεται με το state, όταν το σύστημα φορτώνεται περισσότερο θα αποκτά ταυτόχρονα και μεγαλύτερη ικανότητα αποσυμφόρησης. Άρα, σε σχέση με το μοντέλο σταθερού $\mu$, θα αναμέναμε:

\begin{itemize}
    \item μικρότερη συγκέντρωση πιθανότητας στις υψηλές καταστάσεις,
    \item μικρότερη πιθανότητα blocking,
    \item μικρότερο μέσο αριθμό πελατών $L$,
    \item μικρότερο μέσο χρόνο παραμονής $W$,
    \item και γενικά ταχύτερη επιστροφή του συστήματος από περιοχές υψηλού φόρτου προς χαμηλότερες καταστάσεις.
\end{itemize}

Μάλιστα, επειδή εδώ η βασική τιμή δίνεται ως
\[
\mu = 1 \ \text{πελάτης/sec},
\]
δηλαδή
\[
60 \ \text{πελάτες/min},
\]
οι νέοι ρυθμοί εξυπηρέτησης γίνονται εξαιρετικά μεγάλοι σε σχέση με τις τιμές
\[
\lambda \in \{1,5,10\} \ \text{πελάτες/min}.
\]
Άρα, αν εφαρμοζόταν κατά γράμμα αυτή η νέα παραδοχή, το σύστημα θα γινόταν πολύ λιγότερο φορτισμένο από ό,τι στο προηγούμενο μοντέλο, και οι απώλειες πελατών θα αναμένονταν αισθητά μικρότερες, πιθανότατα σχεδόν αμελητέες στις περισσότερες περιπτώσεις.

\subsection{Συμπέρασμα}

Η τροποποίηση του μοντέλου δεν αλλάζει τη βασική δομή της προσομοίωσης, αλλά αλλάζει ουσιαστικά τον μηχανισμό παραγωγής των departures. Ενώ στην αρχική εκδοχή κάθε χρόνος εξυπηρέτησης παραγόταν από εκθετική κατανομή με σταθερό ρυθμό $\mu$, στη νέα εκδοχή ο ρυθμός εξυπηρέτησης εξαρτάται από την τρέχουσα κατάσταση του συστήματος και γίνεται:
\[
\mu_i=\mu(i+1), \qquad i=1,\ldots,10.
\]

Επομένως, για να προσαρμοστεί σωστά ο κώδικας, πρέπει:

\begin{enumerate}[label=\arabic*.]
    \item να ενοποιηθούν πρώτα οι μονάδες χρόνου,
    \item να αντικατασταθεί ο σταθερός departure ρυθμός από τον state-dependent ρυθμό $\mu_{\text{state}}$,
    \item να αλλάξει ο υπολογισμός του συνολικού ρυθμού και της πιθανότητας επιλογής \\ arrival/departure,
    \item να παραχθεί ο χρόνος μέχρι το επόμενο event με βάση τον νέο συνολικό ρυθμό,
    \item και να επανυπολογίζεται ο μηχανισμός departure μετά από κάθε μεταβολή της κατάστασης.
\end{enumerate}

Όλα τα υπόλοιπα στοιχεία της προσομοίωσης, δηλαδή το blocking, οι counters, οι time-average μετρικές, ο υπολογισμός των $P_i$, $L$, $P_{rej}$ και $W$, καθώς και ο χειρισμός του transient period, παραμένουν ουσιαστικά αμετάβλητα. Συνεπώς, η ζητούμενη επέκταση είναι απολύτως εφικτή πάνω στην υπάρχουσα υλοποίηση, με στοχευμένες αλλαγές μόνο στο κομμάτι που χειρίζεται τον ρυθμό και τον χρόνο της επόμενης αναχώρησης.

\newpage
\section{Παράρτημα Υλοποίησης σε Octave}

\subsection{Κώδικας Υλοποίησης}
\includepdf[pages=-]{lab3.pdf}

\newpage
\subsection{Έξοδος Προγράμματος}
\includepdf[pages=-]{lab3_output.pdf}

\end{document}
